\chapter{Conclusions and Open Mathematical Problems}
\label{ch:ch16_conclusions}

The question of the manuscript was which of the free choices of relativistic field theory the split Clifford algebra $\Cl(3,3)$ and its
tower $\Cl(3n,3n)$ remove, and which they leave free. This chapter states the answer at the strength the proofs support. It names the
strongest results (\cref{sec:contrib}), keeps the mechanisms that failed as results (\cref{sec:failed}), lists the open problems with what
would settle each (\cref{sec:open}), and ends with the test that every statement here has been held to (\cref{sec:testapplied}).

\section{What the manuscript establishes}
\label{sec:contrib}

The following are proved, with the hypotheses stated where they occur. Each is a statement about the algebra or about a specified
construction in it; none is a prediction of a physical model.
\begin{enumerate}[label=(\arabic*),leftmargin=2.2em]
\item \emph{The complex unit and Dirac theory.} Inside the joint commutant of the interior Lorentz action and the internal bivectors, the
only complex structures commuting with the gamma matrices are $\pm J$ (\cref{thm:i}), and the real Dirac theory on $\Cl(3,3)$ is ordinary complex
Dirac theory, exactly (\cref{thm:realification}). The second statement is the reason for the deflationary test of \cref{cor:nopredict}.
\item \emph{Centralisers along the tower.} The bivectors commuting with the interior Lorentz action of $\Cl(3n,3n)$ form $\so{3n-1,3n-3}$,
nested along the tower; at one constituent this is the single generator $J$ (\cref{thm:centraliser,cor:photon1}). The diagonal action instead
gives a compact algebra forming a reductive dual pair with the Lorentz algebra (\cref{thm:casid,cor:howe}).
\item \emph{Two readings of six directions.} Under the centraliser of $G$, $\R^{3,3}$ is $\Lambda^2\R^{1,3}$ with $G$ the Hodge star
(\cref{thm:emfield}), and covariance fixes the form of Maxwell's equations (\cref{cor:maxwellform}). The displacement and field-strength readings
share the metric and the Clifford action, and no equivariant map joins them (\cref{prop:emexclusive,r1:thm:type}).
\item \emph{Interior dynamics.} Factorisation and binding coexist exactly in neutral signature (\cref{thm:neutral}); the force law closes in one
way into a rotor, whose free motion is a pendulum in the duality angle (\cref{thm:closureunique,thm:pendulum}); and the closure has a precise
rigid-body form, whose equations, brackets and conserved quantities are compared with those of the rigid body (\cref{thm:esmap}).
\item \emph{Chiral matter.} The projector $P_*$ of rank $64$ carries one Pati--Salam family with a right-handed neutrino (\cref{thm:PS}); its two
central $\uu1$s cannot be gauged on it (\cref{prop:PSanom}).
\item \emph{Mass kernels.} The Lorentz-invariant mass kernels of the block form a $272$-dimensional space with an explicit complex structure
and Pati--Salam decomposition (\cref{thm:kernelspace}); an adopted potential selects the Standard-Model orbit on a stated open set of couplings
(\cref{thm:select0,thm:compat}).
\item \emph{Charge identification.} At the aligned vacua the electric charge is identified among the compact generators, and the normal charge is
half the fermion number on the chiral block and anomalous (\cref{p3:thm:photon,rem:photonn3}).
\item \emph{Curvature.} The curvature operator of a four-metric lives on the frame, the first Bianchi identity is one trace, and the duality
rotation of the curvature ends at the Einstein tensor, the only divergence-free and variational member of its family; in first-order form the
phase is a Barbero--Immirzi angle (\cref{thm:phase,thm:helmholtz,prop:holst}). These are identities of given metrics.
\item \emph{Causal kinematics.} Choosing a Lorentzian four-plane in $\R^{3,3}$ is choosing a complex structure, and at one constituent the light
cone is the set of chiral spinors squared with the help of $J$ (\cref{thm:selector,thm:cone}).
\item \emph{Several constituents.} Massive constituents can turn at an exactly constant rate on a balance surface (\cref{d6:thm:balance}), and a
field between manifolds acts on frequencies and never on rates (\cref{d7:thm:schr}).
\end{enumerate}

\section{Mechanisms that fail}
\label{sec:failed}

The following mechanisms were tried and fail under the hypotheses stated where each is proved. They are results of the manuscript, not
omissions, and each removes a route that would otherwise look open.
\begin{enumerate}[label=(\arabic*),leftmargin=2.2em]
\item Hypercharge cannot be read off anomaly cancellation for charges built from three commuting involutions (\cref{thm:pairing}).
\item $\su3$ fits in the effective algebra of no axial projector of the five-projector family (\cref{cor:nocolour}).
\item A positive-form interior binds only at zero momentum (\cref{thm:virial}); no Lie--Poisson closure lives on the frame alone
(\cref{thm:noframe}).
\item Non-orthogonality of $\bE$ and $\bB$ generates no mass (\cref{sec:orth:nomass}); momentum along the surplus times lowers the mass
(\cref{en:temporal}); and the $b$-field splits or tears the binding but never creates a rate (\cref{bf:alone,d7:thm:schr}).
\item The normal $U(1)$ of the curved model is not the photon, on four independent counts (\cref{sec:incompatible}), and every normal-plane vector
gravitates with negative energy, in either sign convention (\cref{thm:ghostvec,rem:signconv}).
\item No $SO(3,3)$-invariant scalar order parameter selects the Lorentzian type cleanly: below degree eight it is blind, at every degree the minimum
is degenerate with Lorentz-violating configurations, and its kinetic term has ghosts (\cref{thm:degree,thm:flat,prop:kineticB}).
\item The six-dimensional free field is ill-posed and has no time orientation (\cref{prop:toy,prop:orient}); compact surplus times give
infinitely many tachyons (\cref{ex:tachyons}); Ricci flow does not transfer to split signature (\cref{sec:curved}).
\end{enumerate}

\section{Open problems}
\label{sec:open}

The problems below are sharply posed by the results. For each, the step that would settle it is named.
\begin{enumerate}[label=(\arabic*),leftmargin=2.2em]
\item \emph{Selection among the compact vacua.} A coupling that distinguishes the points of the compact family of vacua at $n\ge3$, where colour
survives, would turn the photon construction into a selection; the present potential encodes its target in its couplings (PH-07).
\item \emph{The electroweak stage.} The selected vacuum is electroweak symmetric and the adopted potential opposes the neutral deformation
(\cref{p3:prop:neutral}); a justified contribution that breaks it, with masses, mixing and the fate of the remaining scalar modes, is missing.
\item \emph{Global minimality.} The global structure of the potential on all $272$ mass kernels, the choice between the Standard-Model vacuum and
its mirror, the electroweak hierarchy, and the explicit breaking of the left-handed fermion number that the axion bound requires are open;
the local statements are reproduced (\cref{thm:fullsecond}), while the global ones rest on archived numerical searches (\cref{prop:randomstart}) and
need a proof.
\item \emph{One theory for both readings.} A finite search over non-diagonal pairings of the two Lorentz bases at the second rung and above would
settle whether the two readings of $\R^{3,3}$ can share one theory at an acceptable price.
\item \emph{The Dirac operator of the blocks on a curved base.} Its spectrum and index are ordinary index-theorem calculations; a nonvanishing
index would not contradict vector-likeness.
\item \emph{The interior and time.} A mechanism placing the interior momenta on the null cone or on the balance surface, the only places where a
constant rate and a bounded frame coexist; the quantisation of the rotor, whose angle would then have to be related to an external time; and
an origin of the scale $\ell$ and of the rate differences between constituents.
\item \emph{Torsion.} Whether the antisymmetric part of the shift is the contorsion of an independent connection is a conjecture that is not
yet type-correct (\cref{sec:transport}).
\item \emph{Gravity.} An action relating metric, pairing, connection and projector, from which the metric follows rather than being given,
with its constraints, degrees of freedom, signs, matter coupling and low-energy limit (\cref{sec:gravlimits}).
\item \emph{Causal order.} A gauge mechanism for type selection, an $SO(3,3)$ analogue of the MacDowell--Mansouri breaking; a principle that
makes the normal translations gauge rather than adding them by hand; constraint preservation for interacting fields; a local clock symmetry
with its connection; a construction of the cone at three constituents without the positive inner product, or a proof that none exists
(\cref{sec:causalverdict}). The universal-evolution hypothesis is open with the requirements of \cref{sec:universaltau}.
\end{enumerate}

\paragraph{Realistic next steps.} One of these is a finite computation that the existing tools can carry: the finite search over non-diagonal
pairings. The index computation is standard. The others are research problems,
and a common local action for a compatible subset of the sectors would be a milestone, not a prerequisite for the mathematical results
stated here.

\section{The test, applied}
\label{sec:testapplied}

The largest gap is the same in every sector, and it is one of selection rather than absence. The gauge algebra is derived and the gauge
field is not; the arena is derived and its dynamics is not; the Einstein tensor of a given metric is constructed and the metric is not
determined; the projector is a field whose equation follows once an action is adopted, but no action is derived and none of its admissible
orbits is singled out; the vacuum of the kernel field is selected, but only on a stated open set of couplings and with the field content
given; and the four-plane on which all of this is built rests on a clock reduction that is imposed. Recovering the hypercharge table, the
anomaly conditions or $\sin^2\theta_W=\tfrac38$ is a consistency result, not evidence, since any construction with that charge content gives
them.

New mathematics would be a theorem that removes one of the choices of \cref{tab:choices} without adding another. New physics would require a
selected spectrum, a coupling or scale relation, a mass restriction, or a deviation that survives re-expression in the complex language of
\cref{thm:realification}. None of these is claimed. What the manuscript does claim is a precise map of where the split-signature Clifford
algebra decides and where it leaves the decision to the physicist: the fixed results, the constructions, the obstructions and the open
problems, each with its hypotheses. On that map the decided territory is larger, and its borders sharper, than the size of the arena would
suggest, and the undecided territory is named.
